% ==============================================================================
% 06_forensics_systeme.tex — UE SIS589
% ==============================================================================
\chapter{Forensique Système~: Disque et Artefacts OS}
\label{chap:forensics-systeme}

\epigraph{\itshape « Les systèmes d'exploitation sont les témoins
  silencieux de nos actions numériques. »}
         {--- Joshua I. James}

\begin{boxobjectifs}
\begin{enumerate}
  \item Analyser les artefacts forensiques Windows~: MFT, Prefetch, registre,
        LNK, shellbags, NTDS.dit.
  \item Analyser les artefacts Linux~: inodes, bashhistory, journal systemd,
        fichiers de configuration.
  \item Reconstruire une timeline NTFS à partir des timestamps MACE.
  \item Récupérer des fichiers supprimés par carving et analyse des inodes.
  \item Utiliser Autopsy et The Sleuth Kit pour l'analyse systématique.
  \item Détecter les marqueurs de persistance (Run keys, services, cron, LD\_PRELOAD).
\end{enumerate}
\end{boxobjectifs}

% ==============================================================================
\section{Forensique Windows}
\label{sec:forensics-windows}
% ==============================================================================

\subsection{Le système de fichiers NTFS~: structures clés}

\begin{boxdef}
\textbf{MFT (\textit{Master File Table})~:} Structure centrale du système
de fichiers NTFS. Chaque fichier ou répertoire est représenté par un
enregistrement MFT de 1~Ko. L'enregistrement contient les attributs du
fichier, dont les timestamps MACE (\textit{Modified, Accessed, Created,
Entry Modified}) et éventuellement les données elles-mêmes (fichiers $<$~700~octets).
La MFT est l'équivalent NTFS de la table des inodes Unix.
\end{boxdef}

\begin{table}[H]
\centering\small
\caption{Timestamps MACE et leur valeur forensique}
\label{tab:mace}
\rowcolors{2}{TableauPair}{TableauImpair}
\begin{tabular}{p{2cm}p{4.5cm}p{7cm}}
\toprule
\tableauHeader{Timestamp} & \tableauHeader{Événement enregistré} &
\tableauHeader{Valeur de détection} \\
\midrule
\textbf{M} (\textit{Modified}) &
Dernière modification du contenu &
Changement de contenu post-infection \\
\textbf{A} (\textit{Accessed}) &
Dernier accès (lecture ou exécution) &
Accès récent malgré une date de création ancienne \\
\textbf{C} (\textit{Changed/Created}) &
Création du fichier &
Comparé à M~: si C~$>$~M, le fichier a été copié \\
\textbf{E} (\textit{Entry Modified}) &
Dernière modification de l'entrée MFT &
Détection du timestomping~: E différent de C/M \\
\bottomrule
\end{tabular}
\end{table}

\begin{boiteRemarque}{Timestomping}
Le timestomping (T1070.006) modifie les timestamps M/A/C via les API Windows
pour que l'attaquant fasse passer son malware pour un fichier légitime ancien.
Mais il ne modifie pas le timestamp E (Entry Modified, stocké dans l'attribut
\texttt{\$STANDARD\_INFORMATION} de la MFT). Un écart entre M/A/C et E
est un indicateur de timestomping.
\end{boiteRemarque}

\begin{lstlisting}[style=StyleTerminal,
  caption={Analyse MFT et timestamps avec MFTECmd et The Sleuth Kit},
  label={lst:mft-analysis}]
# ─── Extraire et analyser la MFT avec MFTECmd (Eric Zimmermann) ─────────────
MFTECmd.exe -f "C:\$MFT" --csv E:\output --csvf mft_export.csv

# Rechercher des fichiers créés pendant la fenêtre d'incident
# dans le CSV, filtrer sur CreatedOn entre 2026-04-12T02:00 et 04:00

# ─── The Sleuth Kit : liste de tous les fichiers (image Linux ext4) ──────────
# fls = file listing depuis une image
fls -r -m "/" srv-web-01-sda.dd > bodyfile.txt
# -r = récursif, -m "/" = format bodyfile (compatible mactime)

# Générer la timeline NTFS avec mactime
mactime -b bodyfile.txt -d 2026-04-12 > timeline-20260412.csv
# -d = filtrer sur une date

# ─── Rechercher des fichiers supprimés (entrées orphelines) ──────────────────
fls -d srv-web-01-sda.dd   # -d = fichiers supprimés uniquement

# Récupérer un fichier supprimé par numéro d'inode
icat srv-web-01-sda.dd 12345 > fichier_recupere.bin
\end{lstlisting}

\subsection{Artefacts Windows clés}

\subsubsection{Registre Windows}

\begin{boxdef}
\textbf{Registre Windows~:} Base de données hiérarchique stockant la
configuration du système et des applications. Cinq ruches principales
(\textit{hives})~: SYSTEM, SOFTWARE, SAM, SECURITY, NTUSER.DAT.
Le registre est une mine forensique~: persistances, MRU, connexions réseau,
logiciels installés, comptes.
\end{boxdef}

\begin{table}[H]
\centering\small
\caption{Clés de registre forensiques prioritaires}
\label{tab:registry-forensics}
\rowcolors{2}{TableauPair}{TableauImpair}
\begin{tabular}{p{8cm}p{5.5cm}}
\toprule
\tableauHeader{Clé de registre} & \tableauHeader{Information forensique} \\
\midrule
\texttt{HKLM\textbackslash Software\textbackslash Microsoft\textbackslash Windows
\textbackslash CurrentVersion\textbackslash Run} &
Programmes démarrés au lancement (persistance T1547.001) \\
\texttt{HKCU\textbackslash Software\textbackslash Microsoft\textbackslash Windows
\textbackslash CurrentVersion\textbackslash Run} &
Persistance par compte utilisateur \\
\texttt{HKLM\textbackslash SYSTEM\textbackslash CurrentControlSet
\textbackslash Services} &
Services installés (T1543.003) \\
\texttt{HKLM\textbackslash SAM\textbackslash SAM\textbackslash Domains
\textbackslash Account\textbackslash Users} &
Comptes locaux, hash NTLM \\
\texttt{HKCU\textbackslash Software\textbackslash Microsoft\textbackslash Windows
\textbackslash CurrentVersion\textbackslash Explorer\textbackslash RecentDocs} &
Fichiers récemment ouverts (MRU) \\
\texttt{HKLM\textbackslash SYSTEM\textbackslash CurrentControlSet
\textbackslash Control\textbackslash TimeZoneInformation} &
Fuseau horaire (critique pour la timeline) \\
\bottomrule
\end{tabular}
\end{table}

\begin{lstlisting}[style=StyleTerminal,
  caption={Analyse de registre avec RegRipper},
  label={lst:registry-forensics}]
# RegRipper : extraction automatique des artefacts de registre
# Depuis l'image montée ou les ruches copiées

# Extraire et analyser la ruche SYSTEM
rip.pl -r /mnt/forensics/Windows/System32/config/SYSTEM \
       -f system > system_report.txt

# Extraire la ruche SOFTWARE (logiciels installés, Run keys)
rip.pl -r /mnt/forensics/Windows/System32/config/SOFTWARE \
       -f software > software_report.txt

# Extraire NTUSER.DAT d'un utilisateur (MRU, historique)
rip.pl -r /mnt/forensics/Users/admin/NTUSER.DAT \
       -f ntuser > ntuser_admin_report.txt

# Clés de Run (persistances) spécifiquement
rip.pl -r /mnt/forensics/Windows/System32/config/SOFTWARE \
       -p run  # -p = plugin spécifique

# Rechercher un nom de malware connu dans le registre exporté
grep -i "lp.sh\|malware\|backdoor" system_report.txt software_report.txt
\end{lstlisting}

\subsubsection{Prefetch}

\begin{boxdef}
\textbf{Prefetch (\texttt{C:\textbackslash Windows\textbackslash Prefetch\textbackslash})~:}
Mécanisme d'optimisation Windows qui enregistre, pour chaque exécutable lancé,
le nombre d'exécutions, la dernière date d'exécution et les DLL chargées.
Existence d'un fichier Prefetch = preuve d'exécution, même si l'exécutable
a été supprimé.
\end{boxdef}

\begin{lstlisting}[style=StyleTerminal,
  caption={Analyse Prefetch avec PECmd (Eric Zimmermann)},
  label={lst:prefetch}]
# PECmd.exe : analyser tous les fichiers Prefetch
PECmd.exe -d "C:\Windows\Prefetch" --csv E:\output --csvf prefetch.csv

# Rechercher l'exécution de fichiers suspects pendant la fenêtre d'incident
# Dans prefetch.csv, filtrer LastRun entre 2026-04-12T02:00 et 04:00

# Exemple de résultat révélateur :
# Executable : WGET.EXE
# LastRun    : 2026-04-12 03:17:03
# RunCount   : 1
# → confirme l'événement Sysmon EID 1 de notre timeline
\end{lstlisting}

\subsubsection{Shellbags et LNK files}

\begin{description}
  \item[\textbf{Shellbags~:}] Clés de registre mémorisant les dossiers
        explorés via l'Explorateur Windows (taille, position, type de vue).
        Forensiquement~: prouvent qu'un utilisateur a navigué dans un
        répertoire, \textbf{même si ce répertoire a été supprimé}.
  \item[\textbf{Fichiers LNK (.lnk)~:}] Raccourcis Windows générés
        automatiquement à l'ouverture d'un fichier. Contiennent le chemin
        complet de la cible, le volume, le MFT record ID, les timestamps.
        Forensiquement~: prouvent l'ouverture d'un fichier (ex.~: document
        exfiltré) même si l'original n'est plus là.
\end{description}

\begin{lstlisting}[style=StyleTerminal,
  caption={Analyse Shellbags et LNK avec LECmd / SBECmd},
  label={lst:shellbags-lnk}]
# LECmd.exe : analyser les fichiers LNK d'un utilisateur
LECmd.exe -d "C:\Users\admin\AppData\Roaming\Microsoft\Windows\Recent" \
           --csv E:\output --csvf lnk_admin.csv

# SBECmd.exe : analyser les shellbags (USRCLASS.DAT)
SBECmd.exe -d "C:\Users\admin\AppData\Local\Microsoft\Windows" \
            --csv E:\output

# Ces outils font partie de la suite Eric Zimmermann Tools (EZTools)
# Téléchargement gratuit sur https://ericzimmermann.github.io/
\end{lstlisting}

\subsubsection{Événements Windows~: journal de sécurité}

Voir chapitre~2 (tableau~\ref{tab:event-ids}). En phase forensique,
l'analyse se fait depuis les fichiers \texttt{.evtx} exportés ou montés
depuis l'image~:

\begin{lstlisting}[style=StyleTerminal,
  caption={Analyse d'un fichier .evtx hors-ligne},
  label={lst:evtx-offline}]
# Lire un fichier .evtx hors-ligne (Python avec python-evtx)
python3 -m evtx /mnt/forensics/Windows/System32/winevt/Logs/Security.evtx | \
  grep -A5 "EventID>4625"

# Avec EvtxECmd (Eric Zimmermann) : extraction CSV
EvtxECmd.exe \
  -f "Security.evtx" \
  --csv E:\output --csvf security_events.csv \
  --inc 4624,4625,4648,4672,4720,4732,1102

# Rechercher les créations de compte sur 24h autour de l'incident
# Dans security_events.csv : EventId==4720, TimeCreated entre 02:00 et 06:00
\end{lstlisting}

% ==============================================================================
\section{Forensique Linux}
\label{sec:forensics-linux}
% ==============================================================================

\subsection{Système de fichiers ext4~: inodes et journalisation}

\begin{boxdef}
\textbf{Inode~:} Structure de données du système de fichiers ext4 contenant
les métadonnées d'un fichier~: propriétaire, permissions, timestamps (atime,
mtime, ctime), pointeurs vers les blocs de données. Le nom du fichier est
stocké dans le répertoire, pas dans l'inode. Supprimer un fichier ne
supprime pas immédiatement son inode ni ses données~: les blocs sont marqués
libres mais leur contenu reste jusqu'à réallocation.
\end{boxdef}

\begin{lstlisting}[style=StyleTerminal,
  caption={Analyse des artefacts Linux depuis une image montée},
  label={lst:linux-forensics}]
# ─── Monter l'image en lecture seule ────────────────────────────────────────
sudo mount -o ro,noatime /dev/loop0p1 /mnt/forensics/

# ─── Informations sur un inode (stat) ────────────────────────────────────────
stat /mnt/forensics/var/log/auth.log
# Affiche : taille, inode, timestamps (Access, Modify, Change), type

# ─── Lister les fichiers par date de modification (récents d'abord) ──────────
find /mnt/forensics/ -newer /mnt/forensics/etc/passwd \
  -type f 2>/dev/null | head -50
# Fichiers modifiés après /etc/passwd (référence temporelle)

# ─── Rechercher des fichiers SUID suspects (élévation de privilèges) ─────────
find /mnt/forensics/ -perm -4000 -type f 2>/dev/null | sort
# SUID légitime : /usr/bin/passwd, /usr/bin/sudo, etc.
# SUID suspect : tout le reste

# ─── Fichiers récemment modifiés dans des répertoires temporaires ─────────────
find /mnt/forensics/tmp /mnt/forensics/var/tmp \
  /mnt/forensics/dev/shm -type f -ls 2>/dev/null

# ─── Rechercher des webshells (PHP dans le DocumentRoot) ─────────────────────
find /mnt/forensics/var/www -name "*.php" \
  -newer /mnt/forensics/etc/apache2/apache2.conf 2>/dev/null

# Analyser le contenu des .php suspects
grep -l "eval\|base64_decode\|system\|exec\|passthru" \
  /mnt/forensics/var/www/html/*.php
\end{lstlisting}

\subsection{Artefacts d'exécution et de persistance Linux}

\begin{table}[H]
\centering\small
\caption{Artefacts de persistance Linux à analyser}
\label{tab:linux-persistance}
\rowcolors{2}{TableauPair}{TableauImpair}
\begin{tabular}{p{5.5cm}p{8cm}}
\toprule
\tableauHeader{Artefact} & \tableauHeader{Analyse forensique} \\
\midrule
\texttt{\textasciitilde/.bash\_history} &
Historique des commandes~; vérifier les commandes suspectes, les téléchargements,
les connexions sortantes \\
\texttt{/etc/cron.d/}, \texttt{/etc/crontab}, \texttt{/var/spool/cron/} &
Tâches planifiées~; tout ajout non documenté est suspect (T1053.003) \\
\texttt{/etc/rc.local}, \texttt{/etc/init.d/} &
Scripts de démarrage (persistances classiques) \\
Unités systemd (\texttt{/etc/systemd/system/*.service}) &
Nouvelles unités créées après la date d'installation du système \\
\texttt{/etc/ld.so.preload} &
LD\_PRELOAD~: injection de bibliothèque au démarrage de tout processus.
Technique de rootkit (T1574.006) \\
\texttt{/proc/[PID]/} &
État des processus en cours (cmdline, maps, net, fd) \\
Clés SSH (\texttt{\textasciitilde/.ssh/authorized\_keys}) &
Clés publiques ajoutées pour persistence SSH (T1098.004) \\
\bottomrule
\end{tabular}
\end{table}

\begin{lstlisting}[style=StyleTerminal,
  caption={Détection des mécanismes de persistance Linux},
  label={lst:linux-persistance}]
# ─── Crontabs de tous les utilisateurs ───────────────────────────────────────
for user in $(cut -f1 -d: /mnt/forensics/etc/passwd); do
  crontab_file="/mnt/forensics/var/spool/cron/crontabs/$user"
  [ -f "$crontab_file" ] && echo "=== $user ===" && cat "$crontab_file"
done
cat /mnt/forensics/etc/crontab
ls -la /mnt/forensics/etc/cron.d/

# ─── Services systemd non standards ─────────────────────────────────────────
find /mnt/forensics/etc/systemd/system/ \
  -name "*.service" -newer /mnt/forensics/etc/os-release | sort

# ─── LD_PRELOAD (rootkit classique) ──────────────────────────────────────────
cat /mnt/forensics/etc/ld.so.preload 2>/dev/null
# Un fichier non vide et non documenté = rootkit probable

# ─── Clés SSH autorisées (backdoors) ─────────────────────────────────────────
find /mnt/forensics -name "authorized_keys" -type f | \
  xargs -I{} sh -c 'echo "=== {} ==="; cat {}'

# ─── Historique bash de tous les utilisateurs ────────────────────────────────
find /mnt/forensics -name ".bash_history" | \
  xargs -I{} sh -c 'echo "=== {} ==="; cat {}' | \
  grep -E "wget|curl|nc |ncat|python|perl|ruby|base64"
\end{lstlisting}

% ==============================================================================
\section{Récupération de fichiers supprimés}
\label{sec:carving}
% ==============================================================================

\begin{boxdef}
\textbf{File Carving~:} Récupération de fichiers supprimés en recherchant
leurs signatures magiques (\textit{magic bytes}) dans l'espace non alloué
ou le slack space, sans s'appuyer sur les métadonnées du système de fichiers.
Fonctionne même si la MFT / les inodes ont été effacés.
\end{boxdef}

\begin{lstlisting}[style=StyleTerminal,
  caption={Récupération de fichiers avec Scalpel et PhotoRec},
  label={lst:file-carving}]
# ─── Scalpel : carving depuis une image ──────────────────────────────────────
# Configurer /etc/scalpel/scalpel.conf (activer les types recherchés)
# Par défaut : pdf, doc, jpg, zip, etc.
scalpel -c /etc/scalpel/scalpel.conf \
        -o /mnt/recovery/scalpel/ \
        srv-web-01-sda.dd

# ─── PhotoRec (testdisk) : carving interactif ────────────────────────────────
photorec /log /d /mnt/recovery/photorec/ srv-web-01-sda.dd

# ─── bulk_extractor : extraction de patterns (emails, URLs, CC) ──────────────
bulk_extractor -R -o /mnt/recovery/bulk_output/ \
               -e net -e email -e url \
               srv-web-01-sda.dd
# Produit des rapports texte de toutes les URLs, emails, IPs trouvés
# dans toute l'image, y compris l'espace non alloué.

# ─── Récupération d'un fichier spécifique par inode (ext4) ───────────────────
# Identifier l'inode du fichier supprimé
debugfs /dev/loop0p1 -R "ls -ld /var/log/auth.log"
# Extraire les blocs par inode
icat -f ext4 srv-web-01-sda.dd 0 987654 > auth.log.recovered
\end{lstlisting}

% ==============================================================================
\section{Autopsy~: analyse systématique}
\label{sec:autopsy}
% ==============================================================================

\begin{boxdef}
\textbf{Autopsy~:} Interface graphique open-source (Sleuth Kit Foundation)
pour l'analyse forensique de disques~: navigation dans les systèmes de
fichiers, récupération de fichiers supprimés, génération de timeline,
carving, extraction de hash, rapport HTML. Standard en formation DFIR.
\end{boxdef}

Modules clés d'Autopsy~:
\begin{itemize}
  \item \textbf{Timeline Analysis~:} visualisation graphique de l'activité
        fichier sur la fenêtre temporelle de l'incident.
  \item \textbf{Keyword Search~:} recherche plein texte dans toute l'image.
  \item \textbf{Hash Lookup~:} comparaison avec NSRL (fichiers système
        légitimes) et bases de malwares connus.
  \item \textbf{Recent Activity~:} extraction des historiques de navigation,
        des documents récents, des téléchargements.
  \item \textbf{Email~:} extraction et indexation des fichiers PST/MBOX.
\end{itemize}

% ==============================================================================
\section*{Résumé}
% ==============================================================================

\begin{boxresume}
\begin{itemize}
  \item NTFS~: MFT, timestamps MACE, détection timestomping (E $\neq$ C/M/A).
  \item Artefacts Windows~: registre (Run keys, SAM), Prefetch (preuves
        d'exécution), LNK (preuves d'ouverture), Shellbags (navigation).
  \item Linux~: inodes ext4, bash\_history, cron, systemd, LD\_PRELOAD,
        authorized\_keys comme vecteurs de persistance.
  \item File carving~: Scalpel, PhotoRec, bulk\_extractor pour récupérer
        données supprimées.
  \item Outils~: The Sleuth Kit (fls, mactime, icat), RegRipper, EZTools
        (MFTECmd, PECmd, LECmd), Autopsy.
  \item Toute analyse sur copie~; jamais sur original~; hash avant/après.
\end{itemize}
\end{boxresume}

\begin{boxexo}
\textbf{Ex.~6.1 — Analyse du registre.}
À partir de l'image \texttt{srv-erp-01.dd} fournie (Lab~2), utilisez
RegRipper pour identifier~: (a)~les clés Run présentes~;
(b)~les comptes locaux dans la SAM~; (c)~le fuseau horaire configuré.
Documentez chaque trouvaille avec sa clé de registre complète.

\medskip\textbf{Ex.~6.2 — Timeline NTFS.}
Avec \texttt{fls} et \texttt{mactime}, générez la timeline complète
de \texttt{srv-web-01.dd} pour la plage 2026-04-12 02:00--06:00 UTC.
Identifiez les 5~événements les plus suspects et annotez-les avec
les TTPs MITRE correspondants.

\medskip\textbf{Ex.~6.3 — Détection de persistance Linux.}
Sur l'image Linux fournie, recherchez et documentez tous les mécanismes
de persistance (cron, systemd, LD\_PRELOAD, authorized\_keys, bashrc).
\end{boxexo}
